\section{Introduction}
\label{sec:intro}

The ion temperature \Ti{} near the top of the pedestal evolves across confinement transitions and through edge-localized-mode cycles~\citep{Ko2015IonTemperature,Viezzer2018IonHeat}, and charge-exchange \Ti{} has served as a feedback quantity on JT-60U~\citep{Yoshida2009RealTime}. Requirement documents ask for the edge ion temperature on time scales of 10 to 100\,ms at an accuracy of about 10\,\%~\citep{Young2008PlasmaMeasurements,Donne2007ChapterDiagnostics}. Monitoring on such time scales is helped by an estimate that exists when a frame is taken. This paper examines one such capability: a causal \emph{nowcast} of the reported carbon-impurity temperature from charge-exchange spectroscopy (CES) on one pedestal-top chord of the KSTAR 2022 campaign, selected by the data provider.

CES obtains \Ti{} from the Doppler width of a line that an impurity ion emits after charge exchange with the atoms of a neutral beam~\citep{Fonck1984DeterminationPlasma}; the KSTAR system observes carbon~\citep{Ko2010ChargeExchange}. In the beam-modulation analysis described for KSTAR, the background of a beam-on frame is estimated from beam-off intervals before and after it~\citep{Lee2020TwoGaussian,Lee2026NeuralNetworkbased}, so the report for a frame need not be complete when the frame is taken. We did not measure when the reports in our files became available, and we assign them no fixed delay. The reports are also conventional labels rather than independent physical truth: edge spectra contain passive emission~\citep{Viezzer2011PassiveEdge}, impurity and main-ion temperatures have been reported to differ at the plasma edge~\citep{Haskey2018MainIon}, and outliers in KSTAR CES data have motivated dedicated detection~\citep{Kim2024KineticProfile}. On the chord studied here some frames carry no usable report, and the reports that exist scatter from frame to frame (\S\ref{sec:data}). The simplest nowcast holds the last report.

One established response works on the spectrum itself. Neural networks have been applied to charge-exchange spectra at JET~\citep{Bishop1993AutomaticAnalysis,Svensson1999AnalysisJet} and more recently on EAST~\citep{Chai2019FastEstimation}, HL-3~\citep{Tian2026RealTimeCxrs} and KSTAR~\citep{Lee2026NeuralNetworkbased}. A simplified fit has been developed for the DIII-D plasma control system~\citep{Chrystal2026RealTime}, and machine-learning fitting has been developed for main-ion spectra~\citep{Haskey2026MachineLearning}. On KSTAR, a two-Gaussian fit~\citep{Lee2020TwoGaussian} and the NN-CES network~\citep{Lee2026NeuralNetworkbased} both address \Ti{} without beam modulation, and ultrafast CES hardware aims to raise the rate of the measurement itself~\citep{Jang2021DevelopmentUltrafast}. These approaches read a spectrum that has been recorded.

Spectrum reading and temporal nowcasting can therefore be viewed as two stages that may be coupled. NN-CES maps a measured CES spectrum to ion temperature and toroidal rotation~\citep{Lee2026NeuralNetworkbased}; its reports could supply the upstream CES history for a temporal estimator. Our downstream estimator uses strictly past CES reports together with current and past BES and ECEI snapshots to estimate the current report before that report is available, and it does not require the current CES spectrum. It estimates a current report; it does not forecast a future physical temperature or profile, as is done in profile prediction from actuators and state history~\citep{Abbate2021DatadrivenProfile}. A faster upstream reader would refresh the anchors on which the nowcast rests, and could equally reduce the marginal need for it; which effect prevails has to be measured. We did not implement this coupling. The present evaluation uses conventional-analysis reports, and coupling to NN-CES remains a possible implementation.

A second body of work estimates slow or missing measurements from other signals. A real-time kinetic reconstruction at DIII-D is designed to operate when charge-exchange input is absent~\citep{Shousha2023MachineLearningbased}. Fast diagnostics have been used to reconstruct slower electron diagnostics at higher time resolution~\citep{Jalalvand2025MultimodalSuperresolution,Imrek2026TemporalSuperresolution}. Missing signals have been reconstructed at HL-3~\citep{Yang2026FusionmaeSelfsupervised}, EAST~\citep{Wang2025TimeSeries} and KSTAR~\citep{Joung2018ImputationFaulty}. Model-based observers merge diagnostics of different rates through a transport model~\citep{Pastore2023ModelBased,Morosohk2023RealTime}. Monotone cubic interpolation~\citep{fritsch1980} and Gaussian-process (GP) regression~\citep{chilenski2015,Kim2024KineticProfile} are standard tools for measured series and profiles; when they fill a frame from the values on both sides of it, they wait for the next value. We do not suggest that this literature is noncausal or restricted to electron quantities. The present study takes the same idea to one measured ion-temperature chord. Its fast inputs are BES~\citep{Lampert2015CombinedHydrogen} and ECEI~\citep{Yun2014Quasi3d} voltage snapshots on the 10\,ms grid. They are used as uncalibrated predictors, not as calibrated \Te{} or density and not as full-rate turbulence measurements; the interpretation of edge ECEI in particular calls for care~\citep{Vanovac2018EffectsDensity}.

The paper is organized around testable comparisons:
\begin{enumerate}\setlength{\itemsep}{1pt}\setlength{\parskip}{0pt}
\item Which causal estimators are useful when scored against hold, a causal GP and two-sided interpolation that waits for the next report (\S\ref{sec:res:best})?
\item What information do BES and ECEI add to the CES history (\S\ref{sec:res:fast})?
\item Do the answers persist under four validation designs, which are given equal weight, and does the computation fit within a 10\,ms frame (\S\ref{sec:res:designs})?
\end{enumerate}
Optional application analyses (\S\ref{sec:res:use}) and the conditions under which the estimator does not help (\S\ref{sec:res:limits}) follow. Two targets are reported. F is the CES report of the current frame. L5 is a level computed offline as the median of at least three of the five subsequent reports. Both are targets only: neither the current F nor the future-derived L5 is an input available online.

The estimators are assembled from familiar models, and no algorithmic priority is claimed; the contribution is an application-specific empirical evaluation. The results concern one chord and one reused campaign. No untouched data exist, so validation intervals should be read with the caution that applies to cross-validation estimates~\citep{Bates2024CrossValidation}. Meeting a computation budget is not a real-time deployment, and no physical mechanism is identified.
